\chapter{Conclusion et Perspectives}

\section{Synthèse des travaux}

Ce mémoire décrit un pipeline Big Data pour la prédiction des maladies cardiovasculaires
à partir du dataset BRFSS 2020--2024.
Le pipeline combine Apache Spark pour le traitement des 2,17 millions d'enregistrements
et deux modèles d'apprentissage automatique, un réseau de neurones profond (DNN)
et XGBoost, comparés sur un protocole d'évaluation unifié.

\subsection{Contributions principales}

\begin{enumerate}
    \item \textbf{Pipeline de binarisation BRFSS :} Un pipeline de transformation spécifique au format BRFSS a été développé, gérant les encodages non standards (1=Oui/2=Non), la variable dérivée \texttt{\_MICHD} pour la cible cardiovasculaire composite, et les cas particuliers de chaque variable (valeurs 88, 7, 9 indiquant des données manquantes ou des refus).

    \item \textbf{Traitement distribué à grande échelle :} Apache Spark a permis l'ingestion et le nettoyage de 2,17 millions d'enregistrements Parquet, incluant la suppression de 416\,535 doublons (19,1\% du dataset brut), avec un calcul distribué de statistiques descriptives.

    \item \textbf{Modèles comparatifs :} Deux modèles ont été entraînés et comparés sur le même jeu de test standardisé :
    \begin{itemize}
        \item DNN (256-128-64, 36 epochs, seuil optimal 0,6713) : AUC-ROC = 0,8191, Recall = 56,3\%
        \item XGBoost (400 arbres, max\_depth=5, learning\_rate=0,05) : AUC-ROC = 0,8195, Recall = 79,9\%
    \end{itemize}

    \item \textbf{Benchmark Spark vs. XGBoost :} Quantification précise de l'avantage de XGBoost standalone (36 à 97 fois plus rapide que Spark GBT en mode local), permettant de déduire que l'apport de Spark dans ce pipeline réside principalement dans le preprocessing distribué et non dans l'entraînement.

    \item \textbf{Dashboard interactif :} Un tableau de bord Streamlit en 3 pages a été développé, permettant la visualisation interactive des résultats sans nécessiter d'infrastructure Big Data.
\end{enumerate}

\subsection{Réponse à la problématique}

La problématique initiale — comment exploiter des données épidémiologiques massives pour prédire le risque cardiovasculaire en comparant ML classique et Deep Learning — a reçu les réponses suivantes :

\begin{itemize}
    \item Spark traite 2,17 millions de lignes Parquet en quelques secondes et supprime
          416\,535 doublons, opération inaccessible à pandas sans segmentation manuelle.
    \item Les class weights ($w_1 = 4{,}847$ pour le DNN, $\textit{scale\_pos\_weight} = 8{,}69$
          pour XGBoost) gèrent le déséquilibre sans créer de données synthétiques,
          produisant un recall de 79,9~\% sur la classe positive avec XGBoost.
    \item Les deux modèles convergent vers le même AUC (0,820), XGBoost étant
          2,6 fois plus rapide à entraîner (67~s vs 178~s) et offrant un meilleur
          recall à seuil identique.
    \item L'absence de HighBP et HighChol dans le fichier pooled~ML plafonne l'AUC
          à 0,82, en deçà des cibles de 0,85 (DNN) et 0,87 (XGBoost).
          Cette contrainte est d'origine documentaire et non algorithmique.
\end{itemize}

\section{Perspectives d'amélioration}

\subsection{Amélioration des données}

\begin{enumerate}
    \item \textbf{Intégration des features manquantes :} Utiliser le fichier EDA BRFSS (qui contient HighBP et HighChol) pour les phases de modélisation. L'ajout de ces deux variables permettrait d'estimer une amélioration de 0,03 à 0,05 en AUC, selon la littérature.

    \item \textbf{Données longitudinales :} Exploiter la dimension temporelle du BRFSS (répondants suivis sur plusieurs années) pour construire des modèles séquentiels (LSTM, Transformer temporel) capturant l'évolution des facteurs de risque.

    \item \textbf{Sources de données complémentaires :} Enrichir avec des données cliniques (ECG, biologie), génomiques ou d'imagerie pour une modélisation multimodale.
\end{enumerate}

\subsection{Amélioration des modèles}

\begin{enumerate}
    \item \textbf{Explicabilité :} Intégrer SHAP (\textit{SHapley Additive exPlanations}) \cite{lundberg2017} pour l'interprétation des prédictions individuelles, essentielle en contexte médical pour la confiance des cliniciens.

    \item \textbf{Calibration des probabilités :} Appliquer une calibration isotonique ou Platt scaling pour améliorer la qualité des probabilités prédites, permettant une meilleure optimisation du seuil de décision.

    \item \textbf{Ensemble de modèles :} Combiner DNN et XGBoost par stacking ou blending pour potentiellement surpasser les performances individuelles.

    \item \textbf{AutoML :} Utiliser des frameworks comme H2O AutoML ou AutoGluon pour explorer automatiquement un espace d'architectures plus large.
\end{enumerate}

\subsection{Amélioration de l'infrastructure}

\begin{enumerate}
    \item \textbf{Cluster Spark distribué :} Évaluer Spark sur un cluster HDFS/YARN réel (8+ nœuds) pour quantifier le speedup en conditions industrielles. À cette échelle, Spark GBT devient compétitif avec XGBoost standalone pour les volumes dépassant 10M de lignes.

    \item \textbf{Déploiement en production :} Créer une API REST (FastAPI ou Flask) exposant le modèle XGBoost comme service de prédiction en temps réel, intégrable dans un système d'information hospitalier.

    \item \textbf{Monitoring MLOps :} Mettre en place un pipeline MLOps (MLflow, DVC) pour le suivi des expériences, la versioning des modèles et la détection du drift de données au fil du temps.
\end{enumerate}

\subsection{Amélioration du dashboard}

\begin{enumerate}
    \item \textbf{Inférence en temps réel :} Intégrer le modèle XGBoost dans le dashboard pour permettre des prédictions individuelles à partir d'un formulaire de saisie des facteurs de risque.
    \item \textbf{Rapport automatisé :} Générer un rapport PDF automatique récapitulant les résultats d'une session d'exploration EDA.
\end{enumerate}

\section{Conclusion générale}

Ce travail montre qu'un pipeline Big Data pour la prédiction des MCV est réalisable
avec des outils open-source (PySpark, TensorFlow, XGBoost, Streamlit)
sur une infrastructure académique à machine unique.

Les performances obtenues, AUC de 0,82 et recall de 79,9~\% pour XGBoost,
sont cohérentes avec la littérature sur des jeux de données BRFSS de composition
comparable. L'écart avec les cibles du PRD tient à des contraintes de données,
pas à des limitations des algorithmes.

Spark est utile pour le preprocessing à grande échelle, pas pour l'entraînement
en mode local. Sur un cluster distribué, la combinaison Spark (ingestion, nettoyage)
et XGBoost (entraînement) reste une architecture valide pour des volumes de l'ordre
de la dizaine de millions de lignes.

Les perspectives les plus directes sont l'intégration de HighBP et HighChol
depuis le fichier BRFSS EDA, l'ajout d'une couche d'explicabilité SHAP,
et le déploiement du modèle XGBoost comme API de prédiction.
